% ============================================================
% Paper 2 - VERSI JICTRA (LaTeX, untuk cek tampilan PDF).
% Konten diturunkan dari paper2-english.tex (versi JISA) + gaya JICTRA (5 bab,
%   referensi nomor [1], narasi sederhana, figur PNG).
% Format JICTRA: English, SI units, abstract 100-200 words, numbered refs,
%   order Title/Authors/Abstract/Body/Acknowledgements/References, <= 20 pages.
%   All NUMBERS come from real experiments (paper2_reviewer_out/*, notebooks).
%   SFM foundation is SUMMARIZED (not repeated from Paper 1) to avoid overlap.
% ============================================================
\documentclass[11pt,a4paper]{article}

\usepackage[utf8]{inputenc}
\usepackage[T1]{fontenc}
\usepackage{amsmath,amssymb}
\usepackage{booktabs}
\usepackage{graphicx}
\graphicspath{{figure-q1/}{figure-p2/}}
\usepackage{geometry}
\usepackage{times}
\usepackage[hidelinks]{hyperref}
\geometry{margin=2.5cm}
\setlength{\parindent}{0pt}
\setlength{\parskip}{0.4em}
\linespread{1.0}

\title{\textbf{Adversarial Robustness of Cross-Network Intrusion Detection
using Semantic Feature Mapping}}

\author{Hero Yudo Martono\thanks{Corresponding author: \texttt{hero@pens.ac.id}},
\quad Iwan Syarif, \quad Ferry Astika Saputra \\[0.3em]
\normalsize Department of Informatics and Computer Engineering, \\
\normalsize Politeknik Elektronika Negeri Surabaya (PENS), Surabaya, Indonesia \\
\normalsize \texttt{hero@pens.ac.id; iwanarif@pens.ac.id; ferryas@pens.ac.id}}

\date{}

\begin{document}
\maketitle

\begin{abstract}
\noindent A network intrusion detector can be robust in one respect yet fragile in
another. Our earlier work showed that a lightweight model over semantically mapped
features, calibrated with only a small amount of target data, generalizes across
different networks. Generalizing across networks, however, does not guarantee
resistance to adversarial evasion, where an attacker slightly alters traffic to avoid
detection. This paper studies how these two properties interact, rather than proposing
a new technique. Using the same nine shared features and a single lightweight model, we
compare four training variants (baseline, few-shot, adversarial training, and
few-shot combined with adversarial training) on two cross-network directions between
CSE-CIC-IDS2018 and UNSW-NB15, under attacks of increasing strength. The results are
reported honestly: cross-network generalization comes from few-shot calibration, while
adversarial training alone does not help it, and combining the two preserves
generalization. The central finding is that robustness on both axes is
direction- and attack-strength-dependent: the same defence can help in one direction
and harm in the other. We further validate the pipeline on real cloud traffic.
\end{abstract}

\medskip
\noindent\textbf{Keywords:} network intrusion detection; adversarial robustness;
cross-network generalization; semantic feature mapping; evasion attack.

% ============================================================
\section{Introduction}
Machine-learning network intrusion detection systems (NIDS) have become a backbone of
modern cyber defence \cite{cremer,khraisat}, and a long line of public datasets and
models has pushed their in-domain accuracy very high. Yet a detector that scores well on
its own benchmark is often fragile along two separate axes. The first is network shift: a
model that excels on one network collapses on another because the traffic distributions
differ \cite{dhooge,cantone}. The second is adversarial evasion: an attacker alters
attack traffic just enough to slip past the detector \cite{biggio2013,goodfellow2015}.
These two weaknesses are usually studied in isolation, but a real deployment must
withstand both at the same time.

Our earlier work addressed the first axis. Using a Semantic Feature Mapping (SFM) that
keeps only the features genuinely comparable across different flow-export tools, together
with a small amount of target-network labels for calibration, a single lightweight
XGBoost recovers cross-network detection and was validated down to real cloud traffic
\cite{martono2024sfm}. That result, however, says nothing about the second axis. A model that generalizes
across networks may still be easy to evade, and it is not obvious whether making it
robust to evasion would damage the cross-network ability it just gained, or whether the
two kinds of robustness can be held at once.

This paper studies exactly that interaction, and it does so as a scientific question
rather than an engineering one. We do not propose a new technique; the building blocks we
use (SFM, few-shot calibration, FGSM and PGD attacks, adversarial training, and a
feature-space constraint that keeps perturbations consistent with protocol arithmetic)
are all established. Our aim is to find out how cross-network generalization and
adversarial robustness behave when they must coexist on one lightweight model, and why
that coexistence turns out to depend on the direction of the network shift. Keeping the
same nine SFM features, the same model, and the same evaluation metric as the earlier
work, we compare four training variants (baseline, few-shot, adversarial training, and
few-shot combined with adversarial training) on the two cross-network directions between
CSE-CIC-IDS2018 and UNSW-NB15, under attacks of increasing strength.

The main finding, not reported before, is that robustness on the two axes is
direction- and attack-strength-dependent: the same defence can help in one direction and
harm in the other. Cross-network generalization comes from the few-shot calibration,
while adversarial training on its own does not close the cross-network gap; adding
adversarial training on top of few-shot preserves generalization, and in one direction it
yields a model robust on both axes at once, yet in the other direction the hardest attack
still breaks every variant. The clearest illustration is that, under the stronger PGD
attack, adversarial training hurts in one direction (MCC falling from 0.524 to 0.302) but
helps in the other (from $-0.124$ to 0.178). We explain this asymmetry through the width
of the decision cells of the tree model on low- versus high-variance target domains, and
we report the trade-offs honestly, including the cases where the defence fails.

The remainder of this paper is organized as follows. Section 2 reviews related work on
cross-network generalization and adversarial robustness in NIDS. Section 3 presents the
proposed method, including the features and model, the four training variants, and the
adversarial generation and threat model. Section 4 reports the performance evaluation
along both axes, the behaviour under increasing attack strength, security metrics beyond
MCC, the mechanism behind the direction dependence, and a validation on real cloud
traffic. Section 5 concludes and outlines future work.

% ============================================================
\section{Literature Review}
This section reviews the two bodies of work that this paper brings together:
cross-network generalization of NIDS and adversarial robustness of NIDS. We keep the
first part short because it is the setting we inherit, and we focus more on the second
part, which is the new axis we study.

Cross-network generalization has been documented repeatedly: supervised detectors that
are excellent on one dataset degrade sharply on another because the traffic
distributions differ \cite{dhooge,verkerken,cantone}. Part of the difficulty is that
different datasets are produced by different flow-export tools, so even features with the
same name are not directly comparable, and benchmark datasets themselves carry labelling
and feature-computation issues \cite{layeghy}. A standard feature set across datasets has
been proposed to ease comparison \cite{sarhan}, and our earlier work went further by
statistically auditing each feature pair and recovering cross-network detection with a
small amount of target-network calibration, using a single lightweight XGBoost over nine
semantically mapped features \cite{xgboost,martono2024sfm}. That earlier work is the
starting point of this paper; we reuse its mapped features and model rather than
re-deriving them.

The second axis, adversarial evasion, has a long history in machine learning. Early work
showed that small, carefully chosen input perturbations can flip a classifier's decision
\cite{biggio2013,szegedy2014,goodfellow2015}, and stronger iterative attacks and careful
evaluation protocols were later introduced to avoid overestimating robustness
\cite{madry2018,carlini2017,athalye2018,carlini2019evaluating}. In the NIDS setting,
however, an attacker cannot change feature values arbitrarily, because network features
are coupled by protocol arithmetic and must correspond to traffic that can actually be
sent. A line of work therefore studies how realistic such attacks are, and how to
constrain perturbations so that they remain consistent with real traffic
\cite{apruzzese2023,merzouk2022,alhajjar2021,sheatsley2022,lin2022idsgan}. We adopt this
view and restrict our attacks to a protocol-consistent feature space rather than claiming
full packet-level deliverability.

A common defence is adversarial training, where adversarial examples are added during
training so that the model learns to resist them \cite{madry2018}. For tree ensembles
such as XGBoost, however, adversarial robustness is less well understood than for neural
networks, and specialized analyses of evasion and hardening of tree models show that the
piecewise-constant decision surface behaves differently from smooth models
\cite{chen2019robusttrees,kantchelian2016}. This matters for our study because our
detector is a tree ensemble.

What is missing from both bodies of work is their intersection. Cross-network studies do
not test adversarial evasion, and adversarial-robustness studies are usually run within a
single dataset, so the question of whether a single model can be robust on both axes at
once, and how the two interact, has not been examined. Table~\ref{tab:p2related}
positions this paper against representative studies along three dimensions: whether
cross-network generalization is addressed, whether adversarial robustness is studied, and
whether the attacks respect feature-space consistency.

\begin{table}[htbp]
\centering
\caption{Positioning against representative work.}
\label{tab:p2related}
\footnotesize
\begin{tabular}{llll}
\toprule
Study & Cross-network & Adversarial & Feature-space constraint \\
\midrule
Cross-dataset studies \cite{dhooge,verkerken,cantone} & Yes & No & --- \\
Standard feature set \cite{sarhan} & Yes & No & --- \\
SFM + few-shot \cite{martono2024sfm} & Yes & No & --- \\
NIDS evasion studies \cite{apruzzese2023,sheatsley2022} & No & Yes & Partial \\
Tree robustness \cite{chen2019robusttrees,kantchelian2016} & No & Yes & No \\
This work & Yes & Yes & Yes \\
\bottomrule
\end{tabular}
\end{table}

% ============================================================
\section{Proposed Method}
We keep the same detection setup as our earlier cross-network work and extend it with an
adversarial evaluation. This section describes the features, data, and model, the four
training variants we compare, and how adversarial examples are generated under
realistic, protocol-consistent constraints.

\subsection{Features, Data, and Model}
The detector is a single lightweight XGBoost \cite{xgboost} trained on the nine shared
features produced by the Semantic Feature Mapping of our earlier work
\cite{martono2024sfm}: duration, forward and backward packet counts, forward and backward
byte counts, forward and backward mean packet sizes, and the source and destination load
(rate) features. We reuse these features as given rather than re-deriving them, so that
this paper can focus on the adversarial axis. Two public datasets act as the two networks,
CSE-CIC-IDS2018 and UNSW-NB15, each reduced to the same nine features and a binary
benign-versus-attack label. Features are standardized with a per-dataset z-score whose
statistics are fit on the training split only, to avoid leakage. We write
A$\rightarrow$B to mean a model trained on source dataset A and tested on target dataset
B, following the usual source-to-target convention, and we evaluate both
CIC$\rightarrow$UNSW and UNSW$\rightarrow$CIC because the two directions are not symmetric.

\subsection{Four Training Variants}
For each direction we compare four training variants. The \emph{baseline} is trained on
the source network only. The \emph{few-shot} variant adds a small amount of
target-network labels (about 1\%) for cross-network calibration. The \emph{adversarial}
variant augments the source data with adversarial examples during training (an
adversarial ratio of 20\% at training strength $\epsilon=0.1$). The proposed
\emph{few-shot + adversarial} variant combines both: it first calibrates with 1\% target
labels and then applies adversarial training on top. Comparing these four isolates what
each ingredient contributes to each axis.

\subsection{Adversarial Generation and Threat Model}
XGBoost is a tree ensemble, so its decision function is piecewise constant and its exact
gradient is zero almost everywhere. We therefore estimate the direction of evasion with a
central finite difference of the loss in z-score space, as in Eq.~\eqref{eq:saliency},
where $\mathcal{L}$ is the binary cross-entropy, $e_i$ the $i$-th feature direction, and
$h=0.01$.

\begin{equation}
S_i(x) = \frac{\mathcal{L}(x + h\,e_i) - \mathcal{L}(x - h\,e_i)}{2h}
\label{eq:saliency}
\end{equation}

A one-step FGSM attack then perturbs the input along the sign of this saliency,
$x_{\mathrm{adv}} = x + \epsilon\,\mathrm{sign}(S(x))$. Because the saliency is computed
from the attacked model itself, the attacker is adaptive and score-based (it needs only
the model's output scores, not the internal tree structure), which gives every variant an
equally hard and realistic test. To avoid overstating robustness with a single weak
attack, we also use a stronger iterative attack (PGD, 10 iterations) and sweep three
attack strengths $\epsilon \in \{0.05, 0.1, 0.2\}$, so that robustness is reported as a
curve rather than a single point. Each experiment is repeated over five seeds
($\{13,42,101,202,303\}$).

Crucially, a network attacker cannot change feature values arbitrarily: packet and byte
counts are non-negative integers, a flow's mean packet size must stay consistent with its
bytes and packets, and an attacker can only add traffic, not remove what was already
sent. We therefore project every perturbation onto a protocol-consistent feature space
(PCFS) that enforces these rules, summarized in Table~\ref{tab:constraint}. This keeps the
attack consistent at the feature level; we do not claim full packet-level deliverability.
Throughout, the primary metric is the Matthews Correlation Coefficient (MCC), which stays
reliable under heavy class imbalance, defined in Eq.~\eqref{eq:mcc} where TP, TN, FP, and
FN are the confusion-matrix counts.

\begin{equation}
\mathrm{MCC} = \frac{\mathrm{TP}\cdot\mathrm{TN} - \mathrm{FP}\cdot\mathrm{FN}}
{\sqrt{(\mathrm{TP}{+}\mathrm{FP})(\mathrm{TP}{+}\mathrm{FN})(\mathrm{TN}{+}\mathrm{FP})(\mathrm{TN}{+}\mathrm{FN})}}
\label{eq:mcc}
\end{equation}

\begin{table}[htbp]
\centering
\caption{Protocol-consistent feature-space (PCFS) constraints.}
\label{tab:constraint}
\footnotesize
\begin{tabular}{lll}
\toprule
Feature & Constraint & Rationale \\
\midrule
all features        & $x \ge 0$                          & no negative counts or rates \\
packet counts       & integer $\ge 0$                    & packets are discrete \\
bytes vs packets    & bytes $\ge$ packets                & each packet carries $\ge 1$ byte \\
mean packet size    & mean $=$ bytes $/$ packets         & keep mean consistent \\
add-only            & increments only ($\Delta \ge 0$)   & attacker may only add traffic \\
load (rate)         & re-derived from bytes $/$ duration & rate consistent with volume \\
\bottomrule
\end{tabular}
\end{table}

% ============================================================
\section{Performance Evaluation}

\subsection{Two-Axis Results: Generalization vs Adversarial Robustness}
% TODO: tabel utama 4 varian x 2 arah (clean target + adaptive evasion), 5-seed.
% Angka nyata: few-shot CIC->UNSW 0.650, UNSW->CIC 0.899; few-shot+adv 0.662/0.898.
% + dua figur PNG (CIC->UNSW, UNSW->CIC).
[ISI]

\subsection{Robustness Under Increasing Attack Strength}
% TODO: FGSM vs PGD across eps {0.05,0.1,0.2}; temuan: FGSM overstate robustness.
% contoh tajam: adv HURTS CIC->UNSW (0.524->0.302) tapi HELPS UNSW->CIC (-0.124->0.178)
[ISI]

\subsection{Security Metrics: Attack Success Rate}
% TODO: ASR (lower better), recall/precision/F1 di bawah adaptive PCFS eps=0.1
[ISI]

\subsection{Why the Interaction is Direction-Dependent}
% TODO: decision-cell width (narasi sederhana) + figur skematik (PNG) + tabel bukti
[ISI]

\subsection{Validation on Real AWS Traffic}
% TODO: 3-tahap (zero-shot -> +1% few-shot -> +adv); recall 0.00->0.97, MCC 0.18-0.25
% field evasion robustness tetap terbuka (jujur)
[ISI]

% ============================================================
\section{Conclusion and Future Work}
% TODO: ringkas temuan dua-sumbu + pesan direction/attack-dependent;
% future work umum (packet-level deliverable attacks, defense yang sadar-arah, dll)
[ISI]

\section*{Acknowledgements}
The authors thank Politeknik Elektronika Negeri Surabaya (PENS) for computing
support.

\section*{Funding}
This research was supported by the Ministry of Higher Education, Science, and Technology
of the Republic of Indonesia.

\section*{Declaration of Generative AI and AI-assisted Technologies}
During the preparation of this work, the authors used AI-assisted tools to help with
English language editing, LaTeX/formatting, restructuring and simplifying the manuscript
text, and generating some of the schematic and plot figures from the authors' own
experimental results. All research design, experiments, data collection, analysis, and
interpretation were performed by the authors, and all reported numbers come from the
authors' own executed, reproducible experiments (code and notebooks are available from
the authors). No results, data, or figures were fabricated by the tools. After using
these tools, the authors reviewed and edited all content and take full responsibility
for the content of the publication.

\begin{thebibliography}{99}
\bibitem{cremer} F.~Cremer, B.~Sheehan, M.~Fortmann, A.~N.~Kia, M.~Mullins, F.~Murphy, and S.~Materne, Cyber risk and cybersecurity: A systematic review of data availability, \emph{Geneva Papers Risk Insur.}, 47(3), pp.~698--736, 2022.
\bibitem{khraisat} A.~Khraisat, I.~Gondal, P.~Vamplew, and J.~Kamruzzaman, Survey of intrusion detection systems: Techniques, datasets and challenges, \emph{Cybersecurity}, 2(1), pp.~1--22, 2019.
\bibitem{biggio2013} B.~Biggio, I.~Corona, D.~Maiorca, B.~Nelson, N.~\v{S}rndi\'c, P.~Laskov, G.~Giacinto, and F.~Roli, Evasion attacks against machine learning at test time, in \emph{Proc.\ ECML PKDD}, pp.~387--402, 2013.
\bibitem{goodfellow2015} I.~J.~Goodfellow, J.~Shlens, and C.~Szegedy, Explaining and harnessing adversarial examples, in \emph{Proc.\ ICLR}, 2015.
\bibitem{martono2024sfm} H.~Y.~Martono, I.~Syarif, and F.~A.~Saputra, Semantic feature mapping for cross-network intrusion detection: validation on a controlled cloud testbed, \emph{under review}, 2024.
\bibitem{moustafa2015} N.~Moustafa and J.~Slay, UNSW-NB15: A comprehensive data set for network intrusion detection systems, in \emph{Proc.\ MilCIS}, pp.~1--6, 2015.
\bibitem{unsweval} N.~Moustafa and J.~Slay, The evaluation of network anomaly detection systems: Statistical analysis of the UNSW-NB15 data set and the comparison with the KDD99 data set, \emph{Inf.\ Secur.\ J.}, 25(1--3), pp.~18--31, 2016.
\bibitem{sharafaldin2018} I.~Sharafaldin, A.~H.~Lashkari, and A.~A.~Ghorbani, Toward generating a new intrusion detection dataset and intrusion traffic characterization, in \emph{Proc.\ ICISSP}, pp.~108--116, 2018.
\bibitem{dhooge} L.~D'Hooge, T.~Wauters, B.~Volckaert, and F.~De~Turck, Inter-dataset generalization strength of supervised machine learning methods for intrusion detection, \emph{J.\ Inf.\ Secur.\ Appl.}, 54, art.~102564, 2020.
\bibitem{verkerken} M.~Verkerken, L.~D'Hooge, T.~Wauters, B.~Volckaert, and F.~De~Turck, Towards model generalization for intrusion detection: Unsupervised machine learning techniques, \emph{J.\ Netw.\ Syst.\ Manage.}, 30(1), pp.~1--25, 2022.
\bibitem{cantone} M.~Cantone, C.~Marrocco, and A.~Bria, Machine learning in network intrusion detection: A cross-dataset generalization study, \emph{IEEE Access}, 12, pp.~144491--144509, 2024.
\bibitem{sarhan} M.~Sarhan, S.~Layeghy, and M.~Portmann, Towards a standard feature set for network intrusion detection system datasets, \emph{Mobile Netw.\ Appl.}, 27(1), pp.~357--370, 2022.
\bibitem{layeghy} S.~Layeghy, M.~Gallagher, and M.~Portmann, Benchmarking the benchmark---Comparing synthetic and real-world network IDS datasets, \emph{J.\ Inf.\ Secur.\ Appl.}, 80, art.~103689, 2024.
\bibitem{xgboost} T.~Chen and C.~Guestrin, XGBoost: A scalable tree boosting system, in \emph{Proc.\ ACM SIGKDD}, pp.~785--794, 2016.
\bibitem{husain} A.~Husain, A.~Salem, C.~Jim, and G.~Dimitoglou, Development of an efficient network intrusion detection model using XGBoost on the UNSW-NB15 dataset, in \emph{Proc.\ IEEE ISSPIT}, pp.~1--7, 2019.
\bibitem{shone} N.~Shone, T.~N.~Ngoc, V.~D.~Phai, and Q.~Shi, A deep learning approach to network intrusion detection, \emph{IEEE Trans.\ Emerg.\ Topics Comput.\ Intell.}, 2(1), pp.~41--50, 2018.
\bibitem{wu} Z.~Wu, H.~Zhang, P.~Wang, and Z.~Sun, RTIDS: A robust transformer-based approach for intrusion detection system, \emph{IEEE Access}, 10, pp.~64375--64387, 2022.
\bibitem{szegedy2014} C.~Szegedy, W.~Zaremba, I.~Sutskever, J.~Bruna, D.~Erhan, I.~Goodfellow, and R.~Fergus, Intriguing properties of neural networks, in \emph{Proc.\ ICLR}, 2014.
\bibitem{madry2018} A.~Madry, A.~Makelov, L.~Schmidt, D.~Tsipras, and A.~Vladu, Towards deep learning models resistant to adversarial attacks, in \emph{Proc.\ ICLR}, 2018.
\bibitem{carlini2017} N.~Carlini and D.~Wagner, Towards evaluating the robustness of neural networks, in \emph{Proc.\ IEEE S\&P}, pp.~39--57, 2017.
\bibitem{athalye2018} A.~Athalye, N.~Carlini, and D.~Wagner, Obfuscated gradients give a false sense of security: Circumventing defenses to adversarial examples, in \emph{Proc.\ ICML}, pp.~274--283, 2018.
\bibitem{carlini2019evaluating} N.~Carlini, A.~Athalye, N.~Papernot, W.~Brendel, J.~Rauber, D.~Tsipras, I.~Goodfellow, A.~Madry, and A.~Kurakin, On evaluating adversarial robustness, \emph{arXiv:1902.06705}, 2019.
\bibitem{apruzzese2023} G.~Apruzzese, M.~Andreolini, L.~Ferretti, M.~Marchetti, and M.~Colajanni, Modeling realistic adversarial attacks against network intrusion detection systems, \emph{ACM Digit.\ Threats Res.\ Pract.}, 3(3), pp.~1--19, 2022.
\bibitem{merzouk2022} M.~A.~Merzouk, F.~Cuppens, N.~Boulahia-Cuppens, and R.~Yaich, Investigating the practicality of adversarial evasion attacks on network intrusion detection, \emph{Ann.\ Telecommun.}, 77, pp.~763--775, 2022.
\bibitem{alhajjar2021} E.~Alhajjar, P.~Maxwell, and N.~Bastian, Adversarial machine learning in network intrusion detection systems, \emph{Expert Syst.\ Appl.}, 186, art.~115782, 2021.
\bibitem{sheatsley2022} R.~Sheatsley, N.~Papernot, M.~Weisman, G.~Verma, and P.~McDaniel, Adversarial examples for network intrusion detection systems, \emph{J.\ Comput.\ Secur.}, 30(5), pp.~727--752, 2022.
\bibitem{lin2022idsgan} Z.~Lin, Y.~Shi, and Z.~Xue, IDSGAN: Generative adversarial networks for attack generation against intrusion detection, in \emph{Proc.\ PAKDD}, pp.~79--91, 2022.
\bibitem{chen2019robusttrees} H.~Chen, H.~Zhang, D.~Boning, and C.-J.~Hsieh, Robust decision trees against adversarial examples, in \emph{Proc.\ ICML}, pp.~1122--1131, 2019.
\bibitem{kantchelian2016} A.~Kantchelian, J.~D.~Tygar, and A.~D.~Joseph, Evasion and hardening of tree ensemble classifiers, in \emph{Proc.\ ICML}, pp.~2387--2396, 2016.
\bibitem{bendavid} S.~Ben-David, J.~Blitzer, K.~Crammer, A.~Kulesza, F.~Pereira, and J.~W.~Vaughan, A theory of learning from different domains, \emph{Mach.\ Learn.}, 79(1--2), pp.~151--175, 2010.
\end{thebibliography}

\end{document}
